
\section{Related Work}
\label{sec:related_work}

\subsection{Video Diffusion and Continuous Control}
Video diffusion and continuous-control methods generate videos from
visual conditions and allow users to adjust specific attributes.
Early models developed spatiotemporal denoising and large-scale video
synthesis~\cite{ho2022video,singer2022make,ho2022imagen}.
Latent video diffusion operates on compressed representations, while
Lumiere and Wan advance spatiotemporal modeling and pretrained video
generation~\cite{blattmann2023align,bar2024lumiere,wan2025wan}.
Motion-control methods guide generation using trajectories, motion
fields, and regional constraints~\cite{yin2023dragnuwa,shi2024motioni2v,
geng2025motionprompting,zhang2025motionpro}.
Go-with-the-Flow uses warped noise, and TokenMotion separates human
and camera motion~\cite{burgert2025gowiththeflow,li2025tokenmotion}.
MotionStone and Mojito control motion intensity, while Text Slider,
FreeSliders, and TokenDial provide fine-grained attribute
control~\cite{shi2025motionstone,he2024mojito,chiu2026textslider,
ezra2025freesliders,liu2026tokendial}.

\subsection{Joint and Function-based Generation}
Joint and function-based generation methods model related outputs
and signals across multiple coordinates.
Multi-view diffusion methods generate consistent images from different
viewpoints: MVDream jointly models camera views, and SyncDreamer uses
relationships among views conditioned on a reference
image~\cite{shi2024mvdream,liu2024syncdreamer}.
Functional generative models represent signals over continuous input
domains. Neural Diffusion Processes model distributions over
functions~\cite{dutordoir2023neural}.
Score-Based Diffusion Models in Function Space and Continuous-Time
Functional Diffusion Processes extend diffusion to function-valued
data~\cite{lim2025score,franzese2023continuous}.
Functional Diffusion studies continuous signal generation, while
HyperDiffusion generates neural-field parameters with
diffusion~\cite{zhang2024functional,erkoc2023hyperdiffusion}.

\subsection{Efficient Diffusion Sampling}
Efficient diffusion sampling methods reduce the computational cost
of generating samples. DDIM introduces an alternative sampling
procedure, while progressive distillation trains models to generate
samples in fewer steps~\cite{song2021denoising,salimans2022progressive}.
Consistency Models and Latent Consistency Models learn consistency
across noise levels to enable fast sampling~\cite{song2023consistency,
luo2023latent}. Other practical strategies for generating related
outputs include reusing the same initial noise to improve consistency
and interpolating the latents of generated endpoints to produce
intermediate results. These approaches support efficient generation
through reduced sampling steps, shared randomness, or interpolation.
